\documentclass[10pt, aspectratio=169]{beamer}
% Preámbulo modular para presentaciones Beamer (Modelación Estadística)
% Diseño y paleta institucional del Tecnológico de Monterrey

\documentclass[aspectratio=169,xcolor={dvipsnames,table}]{beamer}

\usepackage[utf8]{inputenc}
\usepackage[spanish,mexico]{babel}
\usepackage{amsmath,amssymb,amsfonts,mathtools}
\usepackage{graphicx}
\usepackage{listings}
\usepackage{booktabs}
\usepackage{tikz}
\usetikzlibrary{matrix,arrows,decorations.pathmorphing,shapes.geometric,positioning}

% Paleta Institucional Tecnológico de Monterrey
\definecolor{TecAzul}{HTML}{0039A6}       % Azul TEC: color primario institucional
\definecolor{TecAzulOscuro}{HTML}{1A2E51} % Azul marino oscuro
\definecolor{TecRojo}{HTML}{EC2661}       % Rojo de acento (paleta miTec)
\definecolor{TecGris}{HTML}{646464}       % Gris neutro para comentarios y subtítulos
\definecolor{TecFondoBloque}{HTML}{F4F6F9}% Fondo suave para bloques

% Configuración de Tema Beamer
\usetheme{Boadilla}
\usecolortheme[named=TecAzulOscuro]{structure}
\setbeamercolor{alerted text}{fg=TecRojo}
\setbeamercolor{palette primary}{bg=TecAzulOscuro,fg=white}
\setbeamercolor{palette secondary}{bg=TecAzul,fg=white}
\setbeamercolor{palette tertiary}{bg=TecAzulOscuro!80!black,fg=white}
\setbeamercolor{title}{fg=TecAzulOscuro,bg=TecFondoBloque}
\setbeamercolor{frametitle}{fg=TecAzulOscuro,bg=TecFondoBloque}

% Bloques personalizados elegantes
\setbeamercolor{block title}{bg=TecAzulOscuro,fg=white}
\setbeamercolor{block body}{bg=TecFondoBloque,fg=black}
\setbeamercolor{block title alerted}{bg=TecRojo,fg=white}
\setbeamercolor{block body alerted}{bg=TecRojo!8,fg=black}
\setbeamercolor{block title example}{bg=TecAzul,fg=white}
\setbeamercolor{block body example}{bg=TecAzul!8,fg=black}

% Redondear bordes y añadir sombra sutil
\setbeamertemplate{blocks}[rounded][shadow=true]
\setbeamertemplate{navigation symbols}{}

% Pie de página limpio con número de diapositiva
\setbeamertemplate{footline}{
  \leavevmode%
  \hbox{%
  \begin{beamercolorbox}[wd=.7\paperwidth,ht=2.25ex,dp=1ex,left]{author in head/foot}%
    \hspace*{2ex}\insertshortauthor~~(\insertshortinstitute) \hfill \insertshorttitle
  \end{beamercolorbox}%
  \begin{beamercolorbox}[wd=.3\paperwidth,ht=2.25ex,dp=1ex,right]{date in head/foot}%
    \insertshortdate{}\hspace*{2em}
    \insertframenumber{} / \inserttotalframenumber\hspace*{2ex}
  \end{beamercolorbox}}%
  \vskip0pt%
}

% Configuración de Listings de Python para visualización pedagógica de código
\lstset{
    language=Python,
    basicstyle=\ttfamily\footnotesize,
    keywordstyle=\color{TecAzulOscuro}\bfseries,
    stringstyle=\color{TecRojo},
    commentstyle=\color{TecGris}\itshape,
    showstringspaces=false,
    breaklines=true,
    frame=lines,
    rulecolor=\color{TecAzulOscuro!30},
    backgroundcolor=\color{TecFondoBloque},
    aboveskip=0.5em,
    belowskip=0.5em,
    literate={á}{{\'a}}1 {é}{{\'e}}1 {í}{{\'i}}1 {ó}{{\'o}}1 {ú}{{\'u}}1
             {Á}{{\'A}}1 {É}{{\'E}}1 {Í}{{\'I}}1 {Ó}{{\'O}}1 {Ú}{{\'U}}1
             {ñ}{{\~n}}1 {Ñ}{{\~N}}1 {¿}{{?`}}1 {¡}{{!`}}1
}

% Metadatos institucionales por defecto
\title[Modelación Estadística]{Modelación Estadística para Ciencia de Datos e Ingeniería}
\author[J. Castillo Colmenares]{Juliho David Castillo Colmenares}
\institute[Tec de Monterrey]{Tecnológico de Monterrey \\ Escuela de Ingeniería y Ciencias}
\date{\today}

% Comandos y macros estadísticas compartidas para las presentaciones Beamer (Metropolis)

% Conjuntos numéricos básicos
\newcommand{\R}{\mathbb{R}}
\newcommand{\N}{\mathbb{N}}
\newcommand{\Q}{\mathbb{Q}}
\newcommand{\Z}{\mathbb{Z}}
\newcommand{\set}[1]{\left\{ #1 \right\}}
\newcommand{\abs}[1]{\left|#1\right|}
\newcommand{\norm}[1]{\left\| #1 \right\|}

% Comandos estadísticos y probabilísticos del curso
\newcommand{\Var}{\operatorname{Var}}
\newcommand{\cov}{\operatorname{Cov}}
\newcommand{\corr}{\rho}
\newcommand{\comb}[2]{\begin{pmatrix} #1 \\ #2 \end{pmatrix}}
\newcommand{\s}{\sigma}
\newcommand{\particion}{\mathfrak{P}}
\newcommand{\card}[1]{\#\left( #1 \right)}
\newcommand{\seq}[3]{\set{#1}_{#2}^{#3}}
\newcommand{\E}{\mathbb{E}}
\newcommand{\Prob}{\mathbb{P}}

% Entornos pedagógicos y cajas en español académico natural
\newenvironment{discusion}{%
  \begin{alertblock}{Pregunta de Discusión}%
}{%
  \end{alertblock}%
}

\newenvironment{reflexion}{%
  \begin{alertblock}{Reflexión para el Aula}%
}{%
  \end{alertblock}%
}

\newenvironment{paradoja}{%
  \begin{alertblock}{Paradoja de Exploración}%
}{%
  \end{alertblock}%
}

\newenvironment{motivacion}{%
  \begin{exampleblock}{Motivación: Ciencia de Datos en la Práctica}%
}{%
  \end{exampleblock}%
}

\newenvironment{retoclase}{%
  \begin{block}{Reto Rápido en Equipo}%
}{%
  \end{block}%
}


\title[$t$ de Student]{Sección 05.04: Distribución $t$ de Student y Muestras Pequeñas}
\subtitle{Inferencia sobre $\mu$ con $\sigma$ Desconocida}
\author[J. Castillo Colmenares]{Juliho Castillo Colmenares}
\institute{Tecnológico de Monterrey}
\date{\vspace{-1.2cm}}

\begin{document}

% Slide 1: Portada
\begin{frame}[plain]
  \titlepage
\end{frame}

% Slide 2: Hoja de Ruta
\begin{frame}{Hoja de Ruta --- Capítulo 05: Distribuciones de Muestreo}
  \footnotesize
  ¿Dónde nos encontramos en el desarrollo modular del capítulo? \pause
  \begin{itemize}\itemsep=0.08em
    \item \textbf{05.01} Muestreo Aleatorio Simple, Media y Varianza Muestral Insesgada \pause
    \item \textbf{05.02} Teorema del Límite Central Asintótico \pause
    \item \textbf{05.03} Distribución Chi-Cuadrada y Varianza Muestral \pause
    \item \textbf{05.04} Distribución $t$ de Student y Muestras Pequeñas \emph{(Hoy)} \pause
    \item \textbf{05.05} Distribución $F$ de Fisher-Snedecor
  \end{itemize}
  \vspace{-0.05cm}
  \begin{block}{Objetivo de la Sesión}
    Formalizar la distribución $t$ de Student como consecuencia directa del Teorema de Fisher, construir intervalos de confianza exactos para $\mu$ cuando $\sigma$ es desconocida, y entender por qué las muestras pequeñas requieren un ajuste sobre la aproximación Normal.
  \end{block}
\end{frame}

% Slide 3: Motivación
\begin{frame}{El Precio de No Conocer $\sigma$}
  \footnotesize
  \begin{motivacion}
    En la Sección 05.03 vimos que $T=(\bar X-\mu)/(S/\sqrt n) = Z/\sqrt{\chi^2_{n-1}/(n-1)}$: al reemplazar $\sigma$ (desconocida) por $S$ (aleatoria), el estadístico deja de ser exactamente Normal. Hoy formalizamos su distribución exacta y sus consecuencias prácticas.
  \end{motivacion}

  \vspace{0.15cm}
  \pause
  \begin{itemize}\itemsep=0.1cm
    \item \textbf{Colas más pesadas:} $S$ introduce variabilidad adicional respecto a usar $\sigma$ fija. \pause
    \item \textbf{Efecto en $n$ pequeño:} el ajuste es sustancial para muestras chicas, despreciable para grandes. \pause
    \item \textbf{Convergencia:} cuando $n\to\infty$, $t_{n-1}\to N(0,1)$.
  \end{itemize}
\end{frame}

% Slide 4: Definición y Propiedades
\begin{frame}{Definición y Propiedades de la Distribución $t$}
  \footnotesize
  Sean $Z\sim N(0,1)$ y $\chi^2_\nu$ independientes. \pause

  \vspace{0.1cm}
  \begin{block}{Definición}
    $$T = \dfrac{Z}{\sqrt{\chi^2_\nu/\nu}} \sim t_\nu$$
    $$f(t) = \dfrac{\Gamma\left(\frac{\nu+1}{2}\right)}{\sqrt{\nu\pi}\,\Gamma(\nu/2)}\left(1+\dfrac{t^2}{\nu}\right)^{-(\nu+1)/2}, \quad t\in\mathbb{R}$$
  \end{block}

  \vspace{0.1cm}
  \pause
  \begin{alertblock}{Propiedades}
    \begin{itemize}\itemsep=0.05cm
      \item $E(T)=0$ ($\nu>1$); $\Var(T)=\nu/(\nu-2)$ ($\nu>2$) --- siempre mayor que $1=\Var(Z)$.
      \item Colas más pesadas que la Normal para $\nu$ pequeño; $t_\nu\xrightarrow{d}N(0,1)$ cuando $\nu\to\infty$.
    \end{itemize}
  \end{alertblock}
\end{frame}

% Slide 5: Intervalo de Confianza
\begin{frame}{Intervalo de Confianza para $\mu$ con $\sigma$ Desconocida}
  \footnotesize
  Cuando $\sigma$ es desconocida (el caso práctico más común): \pause

  \vspace{0.1cm}
  \begin{block}{Fórmula del Intervalo}
    $$\bar{X} \pm t_{n-1,\,\alpha/2}\cdot\dfrac{S}{\sqrt{n}}$$
  \end{block}

  \vspace{0.1cm}
  \pause
  \begin{alertblock}{¿Por Qué Importa Más en Muestras Pequeñas?}
    Como $t_{n-1,\alpha/2} > z_{\alpha/2}$ siempre, el intervalo con $t$ es más ancho que el que se obtendría (incorrectamente) con $z$. Ejemplo: $t_{9,0.025}\approx 2.262$ vs. $z_{0.025}=1.96$ --- más del $15\%$ de diferencia. Esta diferencia se vuelve despreciable cuando $n\to\infty$.
  \end{alertblock}
\end{frame}

% Slide 6: Prueba t
\begin{frame}{La Prueba $t$ de Una Muestra}
  \footnotesize
  El mismo estadístico permite probar hipótesis sobre $\mu$: \pause

  \vspace{0.1cm}
  \begin{block}{Estadístico de Prueba}
    $$t = \dfrac{\bar{X}-\mu_0}{S/\sqrt{n}} \sim t_{n-1} \text{ bajo } H_0: \mu=\mu_0$$
  \end{block}

  \vspace{0.1cm}
  \pause
  \begin{alertblock}{Decisión}
    Se rechaza $H_0$ a nivel $\alpha$ si $|t| > t_{n-1,\alpha/2}$ (prueba bilateral). Esta es la base de la prueba-$t$ clásica, ampliamente usada cuando $n<30$ y $\sigma$ es desconocida.
  \end{alertblock}
\end{frame}

% Slide 7: Aplicaciones
\begin{frame}{Aplicaciones en Inferencia con Muestras Pequeñas}
  \footnotesize
  La distribución $t$ es la herramienta estándar cuando los datos son costosos o limitados: \pause

  \vspace{0.15cm}
  \begin{itemize}\itemsep=0.12cm
    \item \textbf{Control de calidad:} pruebas de resistencia de materiales con lotes pequeños de producción. \pause
    \item \textbf{Investigación clínica:} ensayos con pocos pacientes; comparaciones antes/después (muestras pareadas). \pause
    \item \textbf{A/B testing temprano:} decisiones con tamaños de muestra iniciales limitados. \pause
    \item \textbf{Preview 05.05:} la distribución $F$ compara varianzas usando el cociente de dos chi-cuadradas, cerrando el trío $\chi^2$--$t$--$F$.
  \end{itemize}
\end{frame}

% Slide 8: Lab Python (Bloque 1)
\begin{frame}[fragile]{Laboratorio en Python: Propiedades y Convergencia de la $t$}
  \scriptsize
  Verificación de $\Var(T)$ y convergencia de cuantiles $t$ a los de la Normal (\texttt{05.04\_student\_t\_distribution.py}):
  {\lstset{basicstyle=\fontsize{4pt}{4.8pt}\selectfont\ttfamily,aboveskip=0.1em,belowskip=0.1em}
  \lstinputlisting[firstline=17, lastline=32]{../../code/05_distribuciones_muestreo/05.04_student_t_distribution.py}}
\end{frame}

% Slide 9: Lab Python (Bloque 2)
\begin{frame}[fragile]{Laboratorio en Python: Intervalos $t$ vs. $Z$}
  \scriptsize
  Comparación de intervalos de confianza basados en $t$ y en $Z$ (incorrecto) para el mismo conjunto de datos:
  {\lstset{basicstyle=\fontsize{4pt}{4.8pt}\selectfont\ttfamily,aboveskip=0.1em,belowskip=0.1em}
  \lstinputlisting[firstline=35, lastline=59]{../../code/05_distribuciones_muestreo/05.04_student_t_distribution.py}}
\end{frame}

% Slide 10: Lab Python (Bloque 3)
\begin{frame}[fragile]{Laboratorio en Python: Prueba $t$ y Cobertura del Intervalo}
  \scriptsize
  Prueba $t$ de una muestra y verificación Monte Carlo de la cobertura del intervalo de confianza:
  {\lstset{basicstyle=\fontsize{4pt}{4.8pt}\selectfont\ttfamily,aboveskip=0.1em,belowskip=0.1em}
  \lstinputlisting[firstline=62, lastline=86]{../../code/05_distribuciones_muestreo/05.04_student_t_distribution.py}}
\end{frame}

% Slide 11: Salida Terminal
\begin{frame}[fragile]{Laboratorio en Python: Salida en Terminal}
  \scriptsize
  Salida estándar generada al ejecutar el script del laboratorio en Python:
  \vspace{0.02cm}
  \begin{block}{Salida Estándar en Consola}
    \fontsize{4pt}{5pt}\selectfont
    \begin{verbatim}
=== Block 1: t-Distribution Properties and Convergence ===
t_10: Var(T)=1.2500 | t_30: Var(T)=1.0714
nu=5: t=2.5706 (+31.15%) | nu=30: t=2.0423 (+4.20%) | nu=120: t=1.9799 (+1.02%)

=== Block 2: t-Based vs Z-Based Confidence Intervals ===
n=9: 95% CI (t): [46.04, 50.96] vs (z, incorrect): [46.41, 50.59] (17.7% wider)
Dataset {23,25,21,27,24}: xbar=24.00, S=2.24, 95% CI: [21.22, 26.78]

=== Block 3: One-Sample t-Test and CI Coverage ===
t=1.75, critical t_24,0.025=2.06 -> Reject H0? No
Monte Carlo CI coverage (n=10): 0.9509 (nominal: 0.9500)
    \end{verbatim}
  \end{block}
\end{frame}

% Slide 12A: Ejercicio Nivel 1 (Enunciado)
\begin{frame}{Ejercicio en Clase --- Nivel 1: Fundamental (1/2)}
  \footnotesize
  \begin{block}{Problema 5.4.1 --- Varianza de la $t$ (Enunciado)}
    Sea $T \sim t_{10}$. Calcule $\Var(T)$.
  \end{block}

  \vspace{0.15cm}
  \pause
  \begin{alertblock}{Planteamiento}
    \begin{itemize}\itemsep=0.04cm
      \item Use directamente $\Var(T)=\nu/(\nu-2)$.
    \end{itemize}
  \end{alertblock}
\end{frame}

% Slide 12B: Ejercicio Nivel 1 (Resolución)
\begin{frame}{Ejercicio en Clase --- Nivel 1: Fundamental (2/2)}
  \scriptsize
  Sustituimos $\nu=10$ directamente: \pause

  \vspace{0.05cm}
  \begin{block}{Cálculo}
    $$\Var(T) = \frac{10}{10-2} = \frac{10}{8} = 1.25.$$
  \end{block}

  \vspace{0.05cm}
  \pause
  \begin{alertblock}{Interpretación}
    \scriptsize
    Como $1.25 > 1 = \Var(Z)$, la $t$ con $\nu=10$ es más dispersa que la Normal estándar, reflejando la incertidumbre adicional de estimar $\sigma$ con $S$. Al aumentar $\nu$, $\Var(T)\to 1$.
  \end{alertblock}
\end{frame}

% Slide 13A: Ejercicio Nivel 2 (Enunciado)
\begin{frame}{Ejercicio en Clase --- Nivel 2: Operativo (1/2)}
  \footnotesize
  \begin{block}{Problema 5.4.5 --- Prueba $t$ de Una Muestra (Enunciado)}
    Se afirma que $\mu_0=100$. Una muestra de $n=25$ produce $\bar X=104.2$ y $S=12$. Calcule $t=(\bar X-\mu_0)/(S/\sqrt n)$ y decida si se rechaza $H_0:\mu=100$ a nivel $\alpha=0.05$, usando $t_{24,0.025}\approx 2.064$.
  \end{block}

  \vspace{0.15cm}
  \pause
  \begin{alertblock}{Estrategia}
    \scriptsize
    Compare $|t|$ contra el valor crítico $t_{24,0.025}$.
  \end{alertblock}
\end{frame}

% Slide 13B: Ejercicio Nivel 2 (Resolución)
\begin{frame}{Ejercicio en Clase --- Nivel 2: Operativo (2/2)}
  \scriptsize
  Calculamos el estadístico: \pause

  \vspace{0.05cm}
  \begin{block}{Cálculo}
    $$t = \frac{104.2-100}{12/5} = \frac{4.2}{2.4} = 1.75$$
  \end{block}

  \vspace{0.05cm}
  \pause
  \begin{alertblock}{Decisión}
    \scriptsize
    Como $|1.75| < 2.064$, \textbf{no se rechaza} $H_0:\mu=100$ a nivel $\alpha=0.05$: la muestra no proporciona evidencia suficiente de que la media difiera de $100$.
  \end{alertblock}
\end{frame}

% Slide 14A: Ejercicio Nivel 3 (Enunciado)
\begin{frame}{Ejercicio en Clase --- Nivel 3: Analítico (1/2)}
  \footnotesize
  \begin{block}{Problema 5.4.7 --- Derivación de $\Var(T)$ (Enunciado)}
    Sabiendo que $E(1/\chi^2_\nu)=1/(\nu-2)$ para $\nu>2$, que $Z$ y $\chi^2_\nu$ son independientes, y que $E(Z^2)=1$, derive $\Var(T)=\nu/(\nu-2)$ para $T=Z/\sqrt{\chi^2_\nu/\nu}$.
  \end{block}

  \vspace{0.1cm}
  \pause
  \begin{alertblock}{Estrategia}
    \scriptsize
    Como $E(T)=0$, $\Var(T)=E(T^2)$. Escriba $T^2$ en términos de $Z^2$ y $\chi^2_\nu$, y use independencia para factorizar la esperanza.
  \end{alertblock}
\end{frame}

% Slide 14B: Ejercicio Nivel 3 (Resolución)
\begin{frame}{Ejercicio en Clase --- Nivel 3: Analítico (2/2)}
  \scriptsize
  Factorizamos usando independencia: \pause

  \vspace{0.05cm}
  \begin{block}{Derivación}
    \scriptsize
    \begin{align*}
      T^2 = \frac{Z^2}{\chi^2_\nu/\nu} = \nu\frac{Z^2}{\chi^2_\nu} \implies
      E(T^2) = \nu\, E(Z^2)\, E\!\left(\frac{1}{\chi^2_\nu}\right) = \nu\cdot 1\cdot\frac{1}{\nu-2} = \frac{\nu}{\nu-2}. \quad \qedhere
    \end{align*}
  \end{block}

  \vspace{0.05cm}
  \pause
  \begin{alertblock}{Interpretación}
    \scriptsize
    Como $E(T)=0$, $\Var(T)=E(T^2)=\nu/(\nu-2)$, confirmando la fórmula. La independencia de $Z$ y $\chi^2_\nu$ es esencial para poder factorizar la esperanza del producto.
  \end{alertblock}
\end{frame}

% Slide 15A: Ejercicio Nivel 4 (Enunciado)
\begin{frame}{Ejercicio en Clase --- Nivel 4: Desafiante (1/2)}
  \footnotesize
  \begin{block}{Problema 5.4.9 --- Muestras Pareadas (Enunciado)}
    En un estudio de dieta, $n=8$ sujetos se pesan antes y después. La diferencia promedio es $\bar d=-3.2$ kg con $S_d=2.1$ kg. Construya un IC del $95\%$ para la diferencia media, usando $t_{7,0.025}\approx 2.365$, y determine si hay evidencia de que la dieta reduce el peso.
  \end{block}

  \vspace{0.15cm}
  \pause
  \begin{alertblock}{Estrategia}
    \scriptsize
    Trate las diferencias $d_i$ como una sola muestra; si el intervalo excluye el cero, hay evidencia de un efecto real.
  \end{alertblock}
\end{frame}

% Slide 15B: Ejercicio Nivel 4 (Resolución)
\begin{frame}{Ejercicio en Clase --- Nivel 4: Desafiante (2/2)}
  \scriptsize
  Calculamos el error estándar y el margen: \pause

  \vspace{0.05cm}
  \begin{block}{Cálculo}
    \scriptsize
    \begin{align*}
      \text{SE} = \frac{2.1}{\sqrt 8} \approx 0.7425, \qquad \text{margen} = 2.365\cdot 0.7425 \approx 1.756.
    \end{align*}
    \begin{align*}
      \text{IC}_{95\%} = -3.2 \pm 1.756 = [-4.956,\ -1.444].
    \end{align*}
  \end{block}

  \vspace{0.05cm}
  \pause
  \begin{alertblock}{Interpretación}
    \scriptsize
    Como el intervalo completo es negativo (excluye el cero), hay evidencia estadística de que la dieta reduce el peso promedio en aproximadamente $1.44$ a $4.96$ kg, con $95\%$ de confianza.
  \end{alertblock}
\end{frame}

% Slide 16: Cierre
\begin{frame}{Cierre de Sección: la $t$ de Student}
  \begin{reflexion}
    La distribución $t$ resuelve el problema práctico más común de la inferencia estadística: estimar $\mu$ cuando $\sigma$ es desconocida. Su construcción --- Normal dividida entre la raíz de una chi-cuadrada independiente --- no es arbitraria, sino la consecuencia directa del Teorema de Fisher que estudiamos la sesión anterior.
  \end{reflexion}

  \vspace{0.2cm}
  \pause
  \begin{block}{Lecciones Clave}
    \begin{itemize}\itemsep=0.08cm
      \item \textbf{Definición:} $T=Z/\sqrt{\chi^2_\nu/\nu}$; colas más pesadas que la Normal, $\Var(T)=\nu/(\nu-2)$.
      \item \textbf{IC para $\mu$:} $\bar X \pm t_{n-1,\alpha/2}\cdot S/\sqrt n$, siempre más ancho que el intervalo basado en $Z$.
      \item \textbf{Convergencia:} $t_\nu \xrightarrow{d} N(0,1)$ cuando $\nu\to\infty$; el efecto es más notorio en muestras pequeñas.
    \end{itemize}
  \end{block}
\end{frame}

% Slide 17: Perspectiva Modular
\begin{frame}{Perspectiva Modular --- El Próximo Reto}
  \scriptsize
  \begin{retoclase}
    Con la $t$ de Student cerramos el par de herramientas para inferencia sobre una sola muestra normal ($\chi^2$ para $\sigma^2$, $t$ para $\mu$). La Sección 05.05 completará el trío con la distribución $F$ de Fisher-Snedecor, construida como el cociente de dos chi-cuadradas independientes, esencial para comparar varianzas entre dos poblaciones y para el Análisis de Varianza (ANOVA).
  \end{retoclase}

  \vspace{0.08cm}
  \pause
  \begin{alertblock}{Preparación Recomendada}
    \begin{itemize}\itemsep=0.06cm
      \item Repasa la propiedad de reproductividad de la chi-cuadrada (Sección 05.03).
      \item Reflexiona sobre por qué el cociente de dos varianzas muestrales independientes sigue una distribución $F$.
    \end{itemize}
  \end{alertblock}
\end{frame}

\end{document}
